\newpage
\thispagestyle{empty}
\null
\vfill
\clearpage

\chapter*{Remerciements}

\large{
À l’approche de la fin de ma troisième année d’alternance, je souhaite exprimer toute ma gratitude envers ceux qui m’ont accompagné dans cette étape déterminante de mon parcours. Cette année est marquée par un renforcement de mes compétences et une immersion toujours plus profonde dans le domaine de la microélectronique.

Je tiens à remercier sincèrement mon tuteur en entreprise, M. Bechetoille, pour son soutien constant et ses conseils précieux. Son accompagnement m’a permis d’aborder cette deuxième année avec confiance, tout en approfondissant ma compréhension des défis et des exigences du métier d’ingénieur.

Je suis également reconnaissant envers mes collègues de l’IP2I, dont l’esprit d’équipe et la bienveillance ont créé un environnement de travail stimulant et encourageant. Leur expertise et leur volonté de partager leurs connaissances ont été une véritable source d’inspiration dans mon développement professionnel.

Je tiens à adresser un remerciement particulier au Dr. Imad Laktineh, responsable du projet PICMIC, grâce à qui j’ai eu l’opportunité de m’impliquer dans un projet de recherche aussi ambitieux et technologiquement avancé. Sa confiance et sa vision m’ont permis de contribuer activement au développement d’un détecteur innovant, au cœur des enjeux actuels en physique des particules. 

Un remerciement particulier est également adressé à M. Labrak, notre responsable de formation PSM, pour sa patience et son soutien indéfectible. Sa capacité à nous guider avec clarté et à proposer des solutions pratiques reste un atout précieux dans notre parcours de formation.

Enfin, je souhaite exprimer ma reconnaissance à mon tuteur académique, M. Massouri, pour son engagement et sa disponibilité, ainsi qu'à l'ensemble des enseignants de la formation PSM et à la direction de CPE Lyon. Leur dévouement et leur passion pour l'enseignement ont été une source de motivation continue et ont contribué de manière significative à mon apprentissage.
}
\clearpage

\newpage
\thispagestyle{empty}
\null
\vfill
\clearpage